\subsection{Softmax classification head (slide 19)}
\label{sec:expl-softmax}

\begin{intuitbox}
We need five scores---one per acuity level---that behave like probabilities.
Then we pick the level with the largest score. That is the entire decision rule.
\end{intuitbox}

\begin{whybox}
Connects the 768-d summary to a clinical class. Say ``largest probability'', not bare \(\arg\max\).
\end{whybox}

Linear map to five logits, then softmax:
\begin{align}
  \mathbf{z} &= W h_{\mathrm{[CLS]}} + \mathbf{b},
  \qquad W\in\mathbb{R}^{5\times 768},\; \mathbf{b}\in\mathbb{R}^{5},
  \label{eq:logits}\\
  \hat{y}_c &= \frac{\exp(z_c)}{\sum_{j=1}^{5}\exp(z_j)},
  \qquad c=1,\ldots,5.
  \label{eq:softmax}
\end{align}
By construction \(\hat{y}_c \ge 0\) and \(\sum_{c=1}^{5}\hat{y}_c = 1\).
\(\hat{y}_c\) is the model's belief that the complaint is acuity level \(c\).

\paragraph{Chosen level.}
\begin{equation}
  \hat{c} = \text{the index \(c\) that maximises \(\hat{y}_c\)}.
\end{equation}

\paragraph{Worked numerical example.}
Illustrative logits for levels \(1\)--\(5\):
\begin{align*}
  \mathbf{z} &= [-1.2,\,-0.5,\,1.8,\,0.3,\,-0.9], \\
  \exp(\mathbf{z}) &\propto [0.30,\,0.61,\,6.05,\,1.35,\,0.41], \\
  \sum_{c} \exp(z_c) &= 8.72, \\
  \hat{\mathbf{y}} &\approx [0.03,\,0.07,\,0.69,\,0.15,\,0.05], \\
  \hat{c} &= 3 \quad\text{(then Yellow via \(f_{\mathrm{SATS}}\))}.
\end{align*}
Deployment logs may round a dominant Yellow near \(72\%\); the mathematics is the same: pick the largest \(\hat{y}_c\).

\begin{center}
\small
\begin{tabular}{@{}lp{9cm}@{}}
\toprule
Symbol & Meaning \\
\midrule
\(h_{\mathrm{[CLS]}}\) & 768-d summary from the final encoder layer \\
\(W,\mathbf{b}\) & Learned classification weights and bias \\
\(z_c\) & Logit (pre-softmax score) for level \(c\) \\
\(\hat{y}_c\) & Probability of level \(c\) \\
\(\hat{c}\) & Predicted level = index of largest \(\hat{y}_c\) \\
\bottomrule
\end{tabular}
\end{center}

\begin{saybox}
``The head turns the summary into five probabilities that add to one. We choose the level with the highest probability.''
\end{saybox}

%----------------------------------------------------------------------
\subsection{Why Stage~3 fine-tuning (loss and one gradient step)}
\label{sec:expl-loss}

\begin{intuitbox}
PubMed pre-training teaches medical language. It does \emph{not} teach five-class SATS acuity on short chief complaints.
Stage~3 fine-tuning on the \(80{,}000\) Triagegeist pairs teaches that mapping.
\end{intuitbox}

\paragraph{Stage 1 in one sentence.}
Masked language modelling on PubMed/PMC teaches \(P(t_i \mid \text{context})\) over biomedical text.
No triage colours yet.

\paragraph{Stage 3 objective (what we need for the talk).}
Given nurse label \(c^\star\) and predicted distribution \(\hat{\mathbf{y}}\), minimise cross-entropy:
\begin{equation}
  \mathcal{L} = -\sum_{c=1}^{5} \mathbf{1}[c=c^\star]\,\log \hat{y}_c = -\log \hat{y}_{c^\star}.
  \label{eq:CE}
\end{equation}
Wrong confident predictions produce large loss.

\paragraph{Backpropagation sketch.}
Gradients flow from the loss through softmax and the linear head into BioBERT:
\begin{equation}
  \frac{\partial \mathcal{L}}{\partial W}
  = \frac{\partial \mathcal{L}}{\partial \hat{\mathbf{y}}}
  \cdot \frac{\partial \hat{\mathbf{y}}}{\partial \mathbf{z}}
  \cdot \frac{\partial \mathbf{z}}{\partial W}.
\end{equation}
In words: (i)~how wrong the probabilities are; (ii)~how softmax transmits that error; (iii)~how logits depend on \(W\).

\paragraph{Gradient descent update.}
\begin{equation}
  W_{\mathrm{new}} = W_{\mathrm{old}} - \alpha \frac{\partial \mathcal{L}}{\partial W},
\end{equation}
with a small learning rate (thesis order \(\alpha = 2\times 10^{-5}\)) so PubMed knowledge is preserved while the head and encoder adapt to urgency labels.

\begin{saybox}
``We fine-tune on eighty thousand labelled complaints with cross-entropy. That is what turns BioBERT language knowledge into five acuity probabilities.''
\end{saybox}

%----------------------------------------------------------------------
\subsection{Colour map \(f_{\mathrm{SATS}}\) and pathway \(P(C)\) (slide 20)}
\label{sec:expl-fsats}

\begin{intuitbox}
Training labels are integers \(1\)--\(5\). The hospital speaks colours, timers, and destinations.
A fixed deterministic map bridges the two languages; the pathway card finishes the product.
\end{intuitbox}

\begin{equation}
  f_{\mathrm{SATS}}(\hat{c}) =
  \begin{cases}
    \mathrm{Red} & \hat{c}=1 \\
    \mathrm{Orange} & \hat{c}=2 \\
    \mathrm{Yellow} & \hat{c}=3 \\
    \mathrm{Green} & \hat{c}\in\{4,5\}
  \end{cases}
  = C_{\mathrm{NLP}}.
  \label{eq:fsats}
\end{equation}
Levels \(4\) and \(5\) both map to Green (routine stream). Blue is a dignity protocol, not an NLP output class here.

\paragraph{Pathway card.}
\begin{equation}
  P(C) = \bigl(T_{\max},\;\mathrm{destination},\;\mathrm{actions},\;\mathrm{escalation}\bigr),
\end{equation}
aligned to \examplehospital{} protocol:
\begin{itemize}
  \item \(T_{\max}\): maximum waiting time for that colour;
  \item destination: primary clinical stream;
  \item actions: support steps (nursing, diagnostics, HIT, \ldots);
  \item escalation: what happens if the timer expires or the patient worsens.
\end{itemize}

\paragraph{Protocol sketch (clinical destinations).}
\begin{center}
\footnotesize
\begin{tabular}{@{}llll@{}}
\toprule
Colour & \(T_{\max}\) & Destination (summary) & Escalation \\
\midrule
Red & 0 min & Resuscitation / Surgery & Immediate senior alert \\
Orange & 10 min & Acute medical / surgical streams & Head nurse if unclaimed \\
Yellow & 60 min & Outpatient urgent / Pediatrics & Re-assess at timer \\
Green & \(<4\) h & Routine OPD / Dental / Eye / PPH & Return if worse \\
\bottomrule
\end{tabular}
\end{center}

\paragraph{Scenario examples (tie to the gate).}
\begin{center}
\footnotesize
\begin{tabular}{@{}lp{9cm}@{}}
\toprule
Scenario & Example phrase \\
\midrule
Gate reject & ``Hello, how are you?'' (no clinical content) \\
Red & ``Crushing chest pain, cannot breathe, sweating heavily'' \\
Orange & ``Severe abdominal pain in pregnancy, bleeding'' \\
Yellow & \runningcomplaint{} \\
Green & ``Routine dental pain for two weeks, no fever'' \\
\bottomrule
\end{tabular}
\end{center}

\begin{saybox}
``Level three becomes Yellow; Yellow becomes a card with about sixty minutes and an urgent destination. That is \(f_{\mathrm{SATS}}\) then \(P(C)\).''
\end{saybox}

%----------------------------------------------------------------------
\subsection{Clinical gate and end-to-end chain (slide 21)}
\label{sec:expl-gate}

\begin{intuitbox}
Do not waste BioBERT on ``hello'' or sports chat.
A cheap rule-based gate asks whether the text looks clinical enough.
\end{intuitbox}

\begin{whybox}
Closes the mathematical story and matches the first box on the pipeline diagram.
\end{whybox}

\paragraph{Relevance score.}
Symptom lexicon hits contribute non-negative scores \(s_j(X)\).
A greeting / non-clinical pattern contributes penalty \(p(X)\) (prototype uses \(p(X)=0.5\) when matched).
\begin{equation}
  g(X) = \mathrm{clip}\!\Bigl(\sum_j s_j(X) - p(X),\, 0,\, 1\Bigr).
  \label{eq:gate}
\end{equation}
Proceed to BioBERT iff \(g(X)\ge 0.35\); otherwise re-prompt for a real chief complaint and assign no colour.
The threshold \(0.35\) rejects obvious non-medical input while accepting short valid complaints such as ``chest pain''.

\paragraph{Worked gate examples.}
Clinical running complaint with hits ``headache'' (\(s_1=0.30\)) and ``feverish'' (\(s_2=0.25\)), no greeting penalty:
\begin{align*}
  \sum_j s_j &= 0.55, \qquad g(X)=\mathrm{clip}(0.55,0,1)=0.55 \ge 0.35
\end{align*}
\(\Rightarrow\) classification proceeds.

Greeting ``Hello, how are you?'' with no symptoms and \(p(X)=0.5\):
\begin{align*}
  g(X)=\mathrm{clip}(0-0.5,0,1)=0 < 0.35
\end{align*}
\(\Rightarrow\) stop; request clinical text.

\paragraph{Full chain (memorise).}
\begin{equation}
  X \xrightarrow{g} \tau(X) \to H^{(0)}
  \xrightarrow{12} h_{\mathrm{[CLS]}}
  \to \hat{\mathbf{y}}
  \xrightarrow{f_{\mathrm{SATS}}} C_{\mathrm{NLP}}
  \to P(C_{\mathrm{NLP}}).
  \label{eq:e2e}
\end{equation}
If \(g(X)<0.35\), the chain stops after the gate.

\begin{enumerate}
  \item Gate: compute \(g(X)\); stop if below \(0.35\).
  \item Tokenise and embed \(\to H^{(0)}\).
  \item Twelve encoder layers \(\to h_{\mathrm{[CLS]}}\).
  \item Softmax head \(\to \hat{\mathbf{y}}\); choose largest component \(\hat{c}\).
  \item Map \(\hat{c}\to C_{\mathrm{NLP}}\) and attach \(P(C_{\mathrm{NLP}})\).
\end{enumerate}

\begin{saybox}
``If the gate says yes, we tokenise, embed, run twelve layers, softmax, map to a colour, and emit a pathway card. That single chain is the mathematical contribution of the thesis.''
\end{saybox}
